\documentclass{article}
\usepackage[utf8]{inputenc}
\usepackage{amsmath, amssymb, amsfonts}
\usepackage{graphicx}
\usepackage{hyperref}
\usepackage[margin=1in]{geometry}

\title{K\"oMaL Level A Math Problems \& Solutions (Sep 2023)}
\author{Auto-generated}
\date{}

\begin{document}
	\maketitle
	
	\section*{Problem 1}
	\noindent\url{https://www.komal.hu/feladat?a=feladat&f=A857&l=en}\\[1em]
	
	\begin{figure}[h!]
		\centering
		\includegraphics[max width=\textwidth]{komal_202309_Level_A_Images/A.gif}
	\end{figure}
	
	A. 857.Let \(\displaystyle ABC\) be a given acute triangle, in which \(\displaystyle BC\) is the longest side. Let \(\displaystyle H\) be the orthocenter of the triangle, and let \(\displaystyle D\) and \(\displaystyle E\) be the feet of the altitudes from \(\displaystyle B\) and \(\displaystyle C\), respectively. Let \(\displaystyle F\) and \(\displaystyle G\) be the midpoints of sides \(\displaystyle AB\) and \(\displaystyle AC\), respectively. \(\displaystyle X\) is the point of intersection of lines \(\displaystyle DF\) and \(\displaystyle EG\). Let \(\displaystyle O_1\) and \(\displaystyle O_2\) be the circumcenters of triangles \(\displaystyle EFX\) and \(\displaystyle DGX\), respectively. Finally, \(\displaystyle M\) is the midpoint of line segment \(\displaystyle O_1O_2\). Prove that points \(\displaystyle X\), \(\displaystyle H\) and \(\displaystyle M\) are collinear.
	
	Proposed byBoldizsár Varga,Verőce
	
	(7 pont)
	
	Deadline expired on October 10, 2023.
	
	\begin{figure}[h!]
		\centering
		\includegraphics[max width=\textwidth]{komal_202309_Level_A_Images/1e5d79e552ba15847ae9e715c.png}
	\end{figure}
	
	We will show that all three points lie on the Euler line of triangle \(\displaystyle ABC\). It is known that \(\displaystyle H\) does, so it is sufficient to prove this for \(\displaystyle M\) and \(\displaystyle X\).
	
	Let \(\displaystyle K\) be the center of the Feuerbach circle of triangle \(\displaystyle ABC\). It is known that \(\displaystyle K\) also lies on the Euler line of triangle \(\displaystyle ABC\). Let \(\displaystyle P\) and \(\displaystyle Q\) be the second intersection points of the Feuerbach circle with segments \(\displaystyle BD\) and \(\displaystyle CE\), respectively. It is known that \(\displaystyle D,\) \(\displaystyle E,\) \(\displaystyle F,\) and \(\displaystyle G\) lie on the Feuerbach circle. Since \(\displaystyle \angle FEQ = 90^{\circ}\), it follows from the Thales' theorem that \(\displaystyle Q\) is the point opposite to \(\displaystyle F\) and similarly, \(\displaystyle P\) is opposite to \(\displaystyle G\) on the Feuerbach circle. Therefore, lines \(\displaystyle FQ\) and \(\displaystyle GP\) pass through point \(\displaystyle K\). Applying Pascal's theorem to the hexagon \(\displaystyle EQFDPG\) inscribed in the Feuerbach circle, we get that lines \(\displaystyle EQ\cap DP = H,\) \(\displaystyle QF\cap PG = K,\) and \(\displaystyle FD\cap GE = X\) are collinear, which means that \(\displaystyle X\) also lies on the Euler line of triangle \(\displaystyle ABC\).
	
	We finish the proof by showing that \(\displaystyle M\) is the midpoint of segment \(\displaystyle XK\). To prove this, we show that quadrilateral \(\displaystyle XO_1KO_2\) is a parallelogram, furthermore, its opposite sides are perpendicular to lines \(\displaystyle AC\) and \(\displaystyle AB\). Due to symmetry, it is sufficient to prove this for lines \(\displaystyle O_2K\) and \(\displaystyle O_1X\). It is clear that \(\displaystyle O_2K\) satisfies this condition since both \(\displaystyle O_2\) and \(\displaystyle K\) lie on the perpendicular bisector of segment \(\displaystyle DG\). Let \(\displaystyle S\) be the intersection of line \(\displaystyle O_1X\) and \(\displaystyle AC\), and let \(\displaystyle \angle DGE = \alpha\). Then, \(\displaystyle \angle DFE = \alpha\) since \(\displaystyle GDEF\) is a cyclic quadrilateral. From the fact that \(\displaystyle \angle XO_1E = 2\alpha\) since it is a central angle in the circumcircle of triangle \(\displaystyle XEF\), we have \(\displaystyle \angle EXO_1 = 90^{\circ}-\alpha\). Hence \(\displaystyle \angle GXS=90^{\circ}-\alpha\), which implies that \(\displaystyle \angle XSG\) is a right angle, as desired.
	
	\vspace{2em}
	\hrule
	\vspace{2em}
	
	\section*{Problem 2}
	\noindent\url{https://www.komal.hu/feladat?a=feladat&f=A858&l=en}\\[1em]
	
	\begin{figure}[h!]
		\centering
		\includegraphics[max width=\textwidth]{komal_202309_Level_A_Images/A.gif}
	\end{figure}
	
	A. 858.Prove that the only integer solution of the following system of equations is \(\displaystyle u=v=x=y=z=0\): \(\displaystyle uv=x^2-5y^2\), \(\displaystyle (u+v)(u+2v)=x^2-5z^2\).
	
	Proposed byBarnabás Szabó,Budapest
	
	(7 pont)
	
	Deadline expired on October 10, 2023.
	
	Suppose for the sake of contradiction that there exists a solution where not all variables are equal to 0. Let's pick one of these solutions where the sum of the absolute values is the smallest possible.
	
	First we show that none of \(\displaystyle u\), \(\displaystyle v\), \(\displaystyle u+v\) and \(\displaystyle u+2v\) can be divisible by 5.
	
	If at least one of \(\displaystyle u\), \(\displaystyle v\), \(\displaystyle u+v\) and  \(\displaystyle u+2v\) is divisible by 5, then \(\displaystyle x\) is divisible by 5. Thus the right hand side of both equations is divisible by 5, and then at least two of \(\displaystyle u\), \(\displaystyle v\), \(\displaystyle u+v\) and  \(\displaystyle u+2v\) is divisible by 5. A quick check shows that in this case both \(\displaystyle u\) and \(\displaystyle v\) must be divisible by 5, and thus the left hand sides of the equations are divisible by 25. This implies \(\displaystyle y\) and \(\displaystyle z\) must also be divisible by 5. However, this makes it possible to divide all variables by 5, and then we get a solution with a smaller sum of the absolute values, which contradicts our assumption. Thus none of \(\displaystyle u\), \(\displaystyle v\), \(\displaystyle u+v\) and  \(\displaystyle u+2v\) is divisible by 5.
	
	The next claim is the key to the solution: if \(\displaystyle p\) is an odd prime that divides \(\displaystyle x^2-5y^2\) and \(\displaystyle p\) does not divide \(\displaystyle x\) and \(\displaystyle y\), then \(\displaystyle p\) is 1 or -1 mod 5. The reason is that this is equivalent to 5 being a quadratic remainder mod \(\displaystyle p\) (being the square of \(\displaystyle x/y\)), and by applying quadratic reciprocity \(\displaystyle p\) is a quadratic remainder mod 5 since 5 is a prime that is 1 mod 4. Thus \(\displaystyle p\) must be 1 or -1 mod 5.
	
	This implies the following: if \(\displaystyle p\) is odd and it is 2 or 3 mod 5, then the exponent of \(\displaystyle p\) in \(\displaystyle x^2-5y^2\) is even. And this means that all powers of primes different from 5 in \(\displaystyle x^2-5y^2\) will be 1 or -1 mod 5. Notice that the same is also true for \(\displaystyle 2\): if \(\displaystyle x\) and \(\displaystyle y\) are both odd, then \(\displaystyle x^2-5y^2\) is 4 mod 8, so the exponent of 2 is 2, and if both are even, we can divide them by 2, and \(\displaystyle x^2-5y^2\) will be divided by 4.
	
	Now if both \(\displaystyle u\) and \(\displaystyle v\) are divisible by \(\displaystyle p\neq 5\), then all variables could be divided by \(\displaystyle p\), which is a contradiction. If only one of them is divisible by \(\displaystyle p\), then according to the above the power of \(\displaystyle p\) in the prime factorization of it will be 1 or -1 mod 5.
	
	If both \(\displaystyle u+v\) and \(\displaystyle u+2v\) are divisible \(\displaystyle p\neq 5\), then both \(\displaystyle u\) and \(\displaystyle v\) are divisible by \(\displaystyle p\), and we have seen that this is not possible. Thus the same is true for \(\displaystyle u+v\) and \(\displaystyle u+2v\) as before: each prime power in their prime factorization has to be 1 or -1 mod 5.
	
	Since none of \(\displaystyle u\), \(\displaystyle v\), \(\displaystyle u+v\) and \(\displaystyle u+2v\) can be divisible by 5, the above means that all of them  has to be 1 or -1 mod \(\displaystyle 5\). However, it is really easy to check that this cannot happen even for \(\displaystyle u\), \(\displaystyle v\) and \(\displaystyle u+v\), and thus we are done.
	
	\vspace{2em}
	\hrule
	\vspace{2em}
	
	\section*{Problem 3}
	\noindent\url{https://www.komal.hu/feladat?a=feladat&f=A859&l=en}\\[1em]
	
	\begin{figure}[h!]
		\centering
		\includegraphics[max width=\textwidth]{komal_202309_Level_A_Images/A.gif}
	\end{figure}
	
	A. 859.Path graph \(\displaystyle U\) is given, and a blindfolded player is standing on one of its vertices. The vertices of \(\displaystyle U\) are labeled with positive integers between \(\displaystyle 1\) and \(\displaystyle n\), not necessarily in the natural order. In each step of the game, the game master tells the player whether he is in a vertex with degree 1 or with degree 2. If he is in a vertex with degree 1, he has to move to its only neighbour, if he is in a vertex with degree 2, he can decide whether he wants to move to the adjacent vertex with the lower or with the higher number. All the information the player has during the game after \(\displaystyle k\) steps are the \(\displaystyle k\) degrees of the vertices he visited and his choice for each step. Is there a strategy for the player with which he can decide in finitely many steps how many vertices the path has?
	
	Proposed byMárton Németh,Budapest
	
	(7 pont)
	
	Deadline expired on October 10, 2023.
	
	Such a strategy does exist.
	
	An algorithm will decide in each step whether to choose the adjacent vertex with the smaller or the bigger label, and will predict whether we will arrive in a vertex with degree 1. We say that the algorithm made a mistake when it predicts a vertex with degree 1 and we actually step in a vertex with degree 2, or we arrive at a vertex with degree 1 without the algorithm predicting it.
	
	Lemma: for every possible labeling of a path with \(\displaystyle n\) vertices there exists an algorithm which makes no mistake in this path, but it will make a mistake in all other paths that are longer than \(\displaystyle n\).
	
	Based on this lemma it is easy to prove our claim. For every \(\displaystyle k\) there exists finitely many labelings of a path with \(\displaystyle k\) vertices. We will apply our algorithm for all possible labelings of the path for \(\displaystyle k=1\), then \(\displaystyle k=2\) and so on. If the path has length \(\displaystyle n\), we will see a mistake in all paths of length \(\displaystyle k<n\), so we know that the path has length at least \(\displaystyle n\). Now we try the algorithm for all possible labelings of the path with \(\displaystyle n\) vertices. The algorithm won't make a mistake when the assumed labeling coincides with the actual labeling. This will show that the path has length at most \(\displaystyle n\). Thus we found out that the length of the path is \(\displaystyle n\). So this strategy will find the length of the path in finitely many steps.
	
	Proof of the lemma:
	
	Let us consider a path with \(\displaystyle n\) vertices and a corresponding labeling and call it \(\displaystyle U\). We will create an algorithm correspondng to \(\displaystyle U\). We assume that we are at the vertex with label 1. We will perform a sequence of steps that takes us to an endpoint of the path. If we end up in an endpoint sooner, or we don't end up in an endpoint, we know that we haven't started from the vertex with label 1. Now we assume that we start at vertex 2, and follow a sequence of steps taking to an endpoint. If the result is not consistent with our assumption, we know that we haven't started at 2. Now we assume we start at 3, and follow the same algorithm, and so on. This is a finite algorithm (it cannot perform more than \(\displaystyle n^2\) steps), and if the actual path is \(\displaystyle U\), we will eventually find an endpoint.
	
	We can forget about the walk we have done so far, and now our aim is to create a finite algorithm that starts from an endpoint and can be performed in \(\displaystyle U\), but will lead to a mistake in every path that is longer than \(\displaystyle U\).
	
	Without loss of generality we can assume that we started from the left endpoint. In the first step we will surely land on the second vertex from the left.
	
	For \(\displaystyle k\le \frac{n+1}{2}\) we know that we stand on the \(\displaystyle k^{th}\) vertex from the left. After this we will perform the sequence of steps that would take us from the \(\displaystyle k^{th}\) vertex from the left to the left endpoint in \(\displaystyle U\), and the algorithm predicts that we are in a vertex with degree 1. This works in \(\displaystyle U\), however, in a path that is longer than \(\displaystyle U\) this will lead to a mistake, or we went straight back to the left endpoint, since \(\displaystyle k\le \frac{n+1}{2}\) guarantees that we cannot end up in the right endpoint in \(\displaystyle k-1\) steps. Now we perform those steps that took us to the \(\displaystyle k^{th}\) vertex from the left endpoint. This will take us there again. Now we check whether at the beginning of this step we moved to the smaller or the bigger neighbour from the \(\displaystyle k^{th}\) vertex. We know that this leads in the direction of the left endpoint, so we will make a move in the opposite direction. This guarantees that we are at the \(\displaystyle (k+1)^{th}\) vertex after this step.
	
	We iterate this until \(\displaystyle k\le \frac{n+1}{2}\).
	
	If \(\displaystyle n\) is even, we know that we will ultimately arrive at the \(\displaystyle \left(\frac{n}{2}+1\right)^{st}\) vertex from the left, and if \(\displaystyle n\) is odd, at the \(\displaystyle \left(\frac{n+3}{2}\right)^{nd}\) vertex. Now we perform the sequence of steps that would take us to the right endpoint from here in \(\displaystyle U\), and predict that we will arrive in a vertex with degree 1. This is correct in \(\displaystyle U\), however in path that is longer than \(\displaystyle U\) we don't perform enough steps to end up in an endpoint, so we will encounter a mistake.
	
	Thus this algorithm that works correctly in \(\displaystyle U\), however, it will make a mistake in longer paths, thus finishing our proof.
	
	\vspace{2em}
	\hrule
	\vspace{2em}
	
\end{document}
